% Draft 1 of the results, the cost model, the sensitivity analysis and the conclusion, 2026-11-05 (log.md 0067).
% Each paragraph opens with its sentence from writing.md, step 3. The record of each number is in the comment beside it.

\section{Results}
\label{sec:results}

\paragraph{First Verdict.}
On the three evaluation streams, the verdict by the rule fixed before the clock is Go at $p = 5$: 6 of 14 scored cells pass, on all three streams, and 4 of the passes are expiry cells.
At $p = 3$ no cell of 12 passes, and the median score is 1.098, so sharing is slower than the best other kernel there.
At $p = 4$ no cell of 12 passes either, with a median score of 1.299.
The guarded star-row kernel P4 takes 0.094 to 0.323 of the time of the BFCB-IG port on every cell.
% results/gate1-M1D05/README.md

\paragraph{Twelve New Streams.}
On twelve streams that no kernel, parameter or reading of this study had touched before, 22 of 36 scored cells pass at $p = 5$, on 11 of the 12 streams.
Pooled with the first set, 28 of 50 cells pass, on 14 of 16 streams, so the verdict holds on the larger set.
The passes fall under all three labels: 9 of 12 insert-only cells, 6 of 12 expiry cells and 7 of 12 real-deletion cells pass on the new streams, so recorded deletions do not prevent a pass.
Figure~\ref{fig:testset} shows every stream with its label, its size and its passes.
% results/gate1-M1D07T/README.md, readings 1 and 2

\paragraph{Batch Size.}
Passes depend on the batch size.
At batches of 1,000 events, 19 of 25 pooled cells pass, and at batches of 100, 9 of 25 pass.
A larger batch holds more updates at each centre, so the row built once at that centre is reused by more spokes.
Figure~\ref{fig:scorebyp} shows the score of every scored cell against the pattern size, for both batch sizes.
% docs/paper/outline.md claim sentence; pass lists of M1D05 and M1D07T

\paragraph{Where It Fails.}
Two streams give no pass: the held-out stream, sx-stackoverflow, and the Wikipedia stream dewiki.
sx-stackoverflow scores 0.846 and 0.978 on its expiry cells, and on its insert-only cells no kernel finished within the limits, so they are not scored.
dewiki, the largest Wikipedia stream, misses the test at 0.710 and 0.783, with medians of 0.669 and 0.671.
% results/gate1-M1D05/README.md, How much weight; results/gate1-M1D07T/README.md, p = 5 table

\paragraph{Triangles and 4-Cycles.}
At $p = 3$, sharing beats the scan by a median of 13\% on the new streams and never reaches the test.
Pooled over 48 scored cells, the median score at $p = 3$ is 0.965, and the best single cell is 0.700.
So the trend that sharing fails for triangles and pays for larger patterns holds only in its second half: for triangles sharing helps a little, and never by the margin the test asks for.
The new streams have no $p = 4$ cells, and on the first three streams no $p = 4$ cell passes.
% results/gate1-M1D07T/README.md, reading 3

\paragraph{What Sharing Beats.}
The kernel that sharing beats is almost always the scan.
At $p = 5$, R2 is the best other kernel on 43 of 50 pooled cells, and the mixed table M5 or the recount A5 on the other 7.
The wedge tables WT and WJ are the best other kernel on no cell.
So the gain measured here is the gain of the shared row over the same kernel without sharing.
% results/gate1-M1D07T/README.md, Best arm at p = 5

\begin{table}[t]
  \centering
  \begin{tabular}{llrrr}
    \hline
    Set & Label & Scored & Pass & Median score \\
    \hline
    First three streams & insert-only and expiry & 14 & 6 & 0.686 \\
    Twelve new streams & real deletions & 12 & 7 & 0.646 \\
                       & synthetic expiry & 12 & 6 & 0.667 \\
                       & insert-only & 12 & 9 & 0.637 \\
    Pooled & all & 50 & 28 & 0.650 \\
    \hline
  \end{tabular}
  \caption{Cells at $p = 5$. A cell passes when the slowest repeat of PR takes at most 0.67 times the fastest repeat of the best other kernel. The score is the ratio of medians. All streams have real timestamps.}
  \label{tab:p5}
\end{table}
% results/gate1-M1D05/README.md; results/gate1-M1D07T/README.md

\section{Why Operation Counts Mislead: a Cost Model}
\label{sec:costmodel}

\paragraph{Two Winners.}
On 60 of the 76 cells of the newer streams, the kernel with the fewest operations is not the fastest kernel.
On 35 of these cells, at $p = 3$, the scan does the fewest operations and the sharing kernel P3 is fastest.
On 25, at $p = 5$, a wedge table, the mixed table or the recount does the fewest operations and PR is fastest.
The flat wedge table WJ runs at a median of 9.55\,ns per counted operation on these streams, against 0.35\,ns for PR.
% results/gate1-M1D07T/README.md, Operations against the clock

\paragraph{Three Models.}
We fit three linear models of a kernel's median time on a cell to its operation counts on that cell, by nonnegative least squares.
Model U charges one price for every operation.
Model K charges one price for each of ten kinds of operation: a list element scanned, a probe, a row read or write, a degree read, a table read, a table write, a graph edit, a step of netting the batch, a step of grouping it, and a step of a static recount.
Model KC adds an extra price for table probes, reads and writes when the kernel's table is larger than the 30\,MB last-level cache.
The models are fitted on one set of streams and judged on another set that no coefficient has seen, by their median relative error and by how often they name the fastest kernel on a cell.
Kernel-cells under 0.05\,s are left out, since their time is dominated by process start and the clock's resolution.
% docs/experiments/m1d09-cost-model.md; scripts/m1d09_cost.py

\paragraph{Fitted Where Parameters Are Chosen.}
Fitted on the three engineering streams and read on the twelve new streams, model KC names the fastest kernel on 63 of 64 cells, at a median error of 23.0\%.
Model U names 21 of 64, at a median error of 71.3\%.
Fitted instead on the first three evaluation streams, model KC names all 64 at 22.9\%.
Figure~\ref{fig:costmodel} shows predicted against measured time for both fits.
% results/cost-M1D10/README.md; results/cost-M1D09/README.md

\paragraph{The Price of a Probe.}
In both fits a table probe costs about 10 to 18\,ns, and a step along a row or an adjacency list about 0.05 to 0.6\,ns.
The operation counts of different kernels are correlated, so a single coefficient is not a measured price of one instruction.
The ratio between the kinds is stable across the two fits, and it is the reason unit cost fails: the table kernels trade many cheap steps for fewer expensive probes.
% results/cost-M1D10/README.md, Coefficients

\paragraph{No Cache-Size Effect.}
The extra prices for tables larger than the cache add nothing in either fit.
On the new streams, model KC names the same 63 winners as model K, and lowers the median error only from 23.5\% to 23.0\%.
The fit on the first three streams gives the same picture.
So the data do not show that a table outgrowing the last-level cache is what makes the table kernels slow.
% results/cost-M1D10/README.md, Cache; results/cost-M1D09/README.md

\paragraph{What the Model Misses.}
The model does not explain the flat table.
It predicts WJ at a median of 0.15 of its measured time with the engineering fit, and 0.18 with the other.
WJ and the tuned table WT have the same operation counts, and WT runs at 1.42\,ns per operation, so what separates them is not in any counter, table size included.
The engineering fit also fails on the first three streams at $p = 3$ and $p = 4$: there it names the wrong winner on 14 of 24 cells, each time preferring the sharing kernel where the scan is fastest, by margins of 0.3\% to 28\%.
At $p = 5$ it names all 12 correctly.
% results/cost-M1D10/README.md, Read beside it; results/cost-M1D09/README.md

\section{Sensitivity}
\label{sec:sensitivity}

\paragraph{The Mixed Table at Its Best Threshold.}
The mixed table was run on the six first passes at every threshold from 16 to 1,024, and the best threshold on each cell was chosen after the fact.
All six passes hold against it.
The best threshold is 16 on the expiry cells and 64 on the rest, and it beats the scan on one stream only, by 9\%.
Figure~\ref{fig:theta} shows the mixed table's time over the threshold against PR and R2.
% results/gate1-M1D06A/README.md

\paragraph{The Scan Margin.}
Most passes rest on a margin of about 10\% over the scan.
If the best other kernel's fastest repeat were 5\% faster on every cell, 22 of 50 pooled cells would pass, and at 10\% faster 16 of 50, below the third the verdict needs.
A faster scan on another machine could therefore change the verdict, and this is the main reason the result is stated for one machine.
% results/gate1-M1D07T/README.md, How much weight the passes carry

\paragraph{Background Load.}
After each timing run, every timed repeat was checked against the recorded start times of the scheduled tasks on the machine.
In the two timing runs of the new streams, overlapping repeats ran a median of 0.00\% and 0.13\% over their kernel's median on the cell, and clean repeats 0.00\%.
Three passes have every repeat of the best other kernel inside a window, and load can only slow that kernel, so for these passes to be false the load would have had to slow it by 15\% to 57\% in every repeat.
Without the two of these passes whose required slowdown is smallest, and one cell that runs in a few milliseconds, 25 of 47 cells pass.
% results/gate1-M1D07T/README.md, Background load and Robustness

\paragraph{An Efficiency Core.}
The 24 cells nearest the pass line were timed again on a performance core and on an efficiency core of the same processor.
On the performance core, each kernel's median is within 1\% of the earlier run, and the pooled count stays at 28 of 50.
On the efficiency core, 24 of 50 cells pass, and every cell lost is a Wikipedia stream with real deletions.
The efficiency core runs PR and R2 about 1.6 times slower, but not by the same factor on every stream.
% results/gate1-M1D08/README.md

\section{Limits and Conclusion}
\label{sec:conclusion}

\paragraph{One Machine.}
All timing comes from one machine.
The efficiency-core run changes the core and its first two cache levels, and keeps the last-level cache and the memory, so it does not replace a second machine.
Given the scan margin of Section~\ref{sec:sensitivity}, a machine with a relatively faster scan could change the verdict at five vertices.
% results/gate1-M1D08/README.md; docs/experiments/m1d08-core-type.md

\paragraph{Sixteen Streams.}
Sixteen real-timestamp streams are few beside the hundreds of instances in static benchmark sets.
Streams with real timestamps are scarce, and those with recorded deletions more so: they are the real instances of the set of Hanauer et al.~\cite{hanauer-dyncount}, chosen by a rule fixed before their sizes were read, and every one of them is used.
We did not add synthetic streams, because a synthetic stream adds instances without adding evidence about real update patterns.
% docs/experiments/m1d07-streams.md; CLAUDE.md hard rule 4

\paragraph{What Is New.}
The mechanism is not new.
DPTL+ and BFCB-IG reuse a row per vertex across a batch~\cite{dptl,bfcb}, and Cheetah groups batches into stars and reports that triangles gain least~\cite{cheetah}.
What this paper adds is the measured gain of that mechanism against the strongest sequential exact baselines, with every cost inside, on real streams with real deletions and a held-out stream, and a cost model that explains why operation counts choose the wrong kernel.

\paragraph{Conclusion.}
Sharing work inside a batch pays for exact counting of 5-cycles against the strongest sequential baselines we could build, on 28 of 50 cells over 14 of 16 real streams, and it does not pay for triangles or 4-cycles by the same test.
It fails on the held-out stream and on the largest Wikipedia stream, and most passes rest on a margin of about 10\% over the scan.
An engineer choosing among these kernels should price a table probe at about twenty or more times a scan step, rather than count operations, because the operation count chooses the wrong kernel on most cells and the priced model chooses the right one on 63 of 64 unseen cells.
